\section*{Figures}

% Figure 1: Gene to Geometry
\begin{figure}[h!]
    \centering
    \includegraphics[width=0.9\textwidth]{fig_gene_to_geometry.pdf}
    \caption{\textbf{Counter-Curvature as Active Morphogenesis.} (A) \textbf{Passive Gravitational Prior}: In the absence of active correction, gravity dictates a low-energy, monotonic C-shaped sag (blue curve). (B) \textbf{Active Counter-Curvature Posterior}: Biological growth, guided by the ``Genetic Prior'' (HOX genes), generates an S-shaped ``Counter-Curvature'' (orange curve) that actively opposes the gravitational metric. This high-energy state is maintained by the continuous expenditure of metabolic work ($P_{\mathrm{counter}}$). (C) \textbf{The Energy Deficit Window}: During rapid adolescent growth, the thermodynamic cost of maintaining this active posterior (scaling as $L^3$) outpaces the diffusive nutrient supply ($L^2$). This triggers a ``Vector-Scalar Mismatch'' where high-cost Vector Sensors (Piezo2) fail, leading to a collapse into the passive, gravity-dominated scoliotic mode.}
    \label{fig:iec_landscape}
\end{figure}

% Figure 2: Mode Spectrum
\begin{figure}[h!]
    \centering
    \includegraphics[width=0.9\textwidth]{fig_mode_spectrum.pdf}
    \caption{\textbf{Gravity-Selected vs. Information-Selected Modes.} (A) First three eigenmodes of a passive uniform beam in gravity; the ground state (Mode 0, blue) is a monotonic C-shaped sag. (B) Eigenmodes of the IEC-coupled beam ($\chi_\kappa=0.05$); the shift in effective stiffness $B_{\mathrm{eff}}(s)$ reorders the spectrum, selecting the S-shaped mode as the new energetic ground state. (C) Normalized eigenvalue spectrum $\lambda_n/\lambda_0$ showing the shift in the lowest energy mode.}
    \label{fig:mode_spectrum}
\end{figure}

% Figure 3: 3D Solutions (Main Panel)
\begin{figure}[h!]
    \centering
    \includegraphics[width=0.9\textwidth]{fig_countercurvature_metrics.png}

    \caption{\textbf{Information-driven countercurvature as Active Inference.}
    (A) Curvature profiles showing the transition from the \textbf{Passive Gravitational Prior} (blue, monotonic sag) to the \textbf{Active Posterior} (orange, S-shape). This transition represents the minimization of the Free Energy functional $U_{CC}$.
    (B) The countercurvature metric $g_{\mathrm{eff}}(s)$ highlights regions where the Bio-Gravitational Number ($\mathcal{B}_g$) is maximized to resist deformation.
    (C) Geodesic deviation $\widehat{D}_{\mathrm{geo}}$ (representing the prediction error $\mathcal{E}$) scales with coupling strength; stability requires the gain $\chi_\kappa$ to exceed a critical threshold.
    (D) Microgravity adaptation ($g \to 0$) and \textbf{Spinal Jetlag}. The system relaxes towards the information reference state. The reduced mechanical strain $\gEff$ shown here corresponds to the \textbf{Convective Shutdown}~\cite{sato2025convective}, where loss of daily loading cycles halts solute transport. This triggers the ``Nuclear Stiffness Gauge'' downregulation (Active Inference failure) and stabilizes HIF-1$\alpha$, locking the spine in a ``permanent night'' swelling state that degrades the matrix ($B_{\mathrm{eff}}$) and precipitates scoliotic buckling.}
    \label{fig:countercurvature_main}
\end{figure}

% Figure 4: Phase Diagram
\begin{figure}[h!]
    \centering
    \includegraphics[width=1.0\textwidth]{phase_diagram.png}
    \caption{\textbf{Phase Diagram of Countercurvature Regimes.} Heatmap of normalized geodesic deviation $\widehat{D}_{\mathrm{geo}}$ in the $(\chi_{\kappa},g)$ plane. Three distinct regimes are identified: (I) Gravity-dominated (sag), (II) Cooperative (stable S-curves matching biological norms), and (III) Information-dominated (high $\chi_\kappa$, leading to metric distortions and potential instability).}
    \label{fig:phase_diagram}
\end{figure}

% Figure 5: Scoliosis Emergence
\begin{figure}[h!]
    \centering
    \includegraphics[width=1.0\textwidth]{fig_scoliosis_emergence.png}
    \caption{\textbf{Scoliosis as Countercurvature Failure.} (A) Lateral scoliosis index $S_{\mathrm{lat}}$ vs coupling strength $\chi_\kappa$ for varying asymmetry levels $\varepsilon$. (B) Bifurcation diagram showing the emergence of large lateral deviations from small asymmetries in the information-dominated regime. (C) Scoliosis regime map in $(\chi_\kappa, \varepsilon)$ space; the shaded region indicates scoliotic configurations (Cobb $> 10^\circ$). (D) Amplification factor ($S_{\mathrm{lat}}/\varepsilon$), showing that in the high-$\chi_\kappa$ regime, small patterning errors are amplified by over 100-fold. This instability models the 'Frustrated Repair' mechanism driven by \textbf{Vector-Scalar Mismatch}. When gravity-dependent Vector Sensors (e.g., Piezo2, Lamin A/C) are silenced by unloading or mutation~\cite{steger2025gravity, touchstone2024nuclear}, the system loses its directional reference. However, Scalar Sensors (e.g., Piezo1) remain active, detecting localized stress concentrations. The system responds by ramping up ``gain'' ($\chi_\kappa$) to correct an error it cannot spatially resolve, amplifying small asymmetries into macroscopic deformities~\cite{djebar2024astrogliosis}. This high-gain instability is consistent with the loss of voltage-dependent gain control in Piezo2 sensors~\cite{sanchez2024piezo}. Furthermore, the 'Inflamed Torsion' pathway—driven by unchecked MMP1/3 activity in the absence of gravitational loading~\cite{chen2025mmp}—degrades torsional stiffness ($GJ$), effectively lowering the critical buckling threshold. This instability is mechanistically underpinned by the Piezo1-Cilia Stress Lock, where compressive stress drives osteogenic differentiation via the Piezo1-Runx2 axis~\cite{chen2025piezo_cilia}, permanently locking the spine into this scoliotic state.}
    \label{fig:scoliosis_emergence}
\end{figure}

% Figure 6: Energy Deficit Window
\begin{figure}[h!]
    \centering
    \includegraphics[width=0.8\textwidth]{energy_deficit_window.png}
    \caption{\textbf{Thermodynamic Collapse: The Energy Deficit Window.} (A) While the gravitational moment scales as $M \propto L^4$ (The Gravity Paradox), the metabolic power $P_{\mathrm{counter}}$ (red) required to maintain the Counter-Curvature active standing wave scales as $L^2$ (assuming isometric growth). (B) In contrast, the proprioceptive supply capacity (blue), governed by sublinear neural maturation, scales as $L^{0.5}$. The intersection at $L_{\mathrm{crit}} \approx 0.35$~m marks the onset of the \textbf{Energy Deficit Window} ($L > L_{\mathrm{crit}}$), resulting in a 41.3\% metabolic deficit at $L=0.45$~m. Structural analysis of \textbf{Thermodynamic Cost Proteins} (Table 2) reveals the molecular basis of this deficit: high-anisotropy sensors like Piezo2 and Lamin A/C (demand-side) are structurally expensive to maintain~\cite{lynch2015bioenergetic}, while compact supply-side enzymes (PLOD1, MMPs) are cheaper but rate-limited. In the deficit regime, the system cannot afford the high-fidelity state, triggering a \textbf{Convective Shutdown}~\cite{sato2025convective} and a collapse into the low-energy scoliotic mode.}
    \label{fig:energy_deficit_window} % Verified
\end{figure}

% Figure 7: Energy Deficit Bifurcation
\begin{figure}[h!]
    \centering
    \includegraphics[width=0.9\textwidth]{energy_deficit_bifurcation.png}
    \caption{\textbf{Bifurcation and Stability Transition in the Energy Deficit Regime.} (A) Equilibrium Cobb angle as a function of the Thermodynamic Instability Ratio $\InstabilityR$. For $\InstabilityR < 1$ (cooperative regime), the sagittal S-curve is stable against lateral perturbations. As $\InstabilityR$ crosses unity, the system undergoes a supercritical pitchfork bifurcation, where the sagittal mode becomes unstable and the scoliotic mode emerges as the global energetic minimum. (B) Phase-space trajectories of the spinal geometry $G(s)$ showing the attractor transition from the Information-Geodesic (S-shape) to the Torsional attractor. This transition quantitatively predicts the critical Cobb angle spike observed in clinical datasets upon crossing the volumetric metabolic threshold.}
    \label{fig:energy_bifurcation}
\end{figure}

% Figure 8: AlphaFold Morphology Landscape
\begin{figure}[h!]
    \centering
    \includegraphics[width=0.9\textwidth]{morphology_space_afcc.png}
    \caption{\textbf{Structural Hierarchy of Thermodynamic Cost Proteins.} Morphology landscape mapping key proteins involved in proprioception ($\eta_p$), active maintenance ($\eta_a$), and basal metabolism ($\Gamma_m$) based on their AlphaFold-predicted anisotropy (extension) and radius of gyration (size). Proteins in the high-anisotropy regime (e.g., Piezo2, Vimentin, Lamin A/C) represent the primary ``Thermodynamic Tax'' on the counter-curvature standing wave. Their extended conformations are structurally expensive to maintain, making them the first points of failure when the system enters the Energy Deficit Window ($L > L_{\mathrm{crit}}$). In contrast, supply-side enzymes (Sirt1, SOX9) cluster in the globular, low-anisotropy regime, creating a molecular mismatch between sensory demand and metabolic supply.}
    \label{fig:morphology_landscape}
\end{figure}
